\section{Related work}
Recent work has surveyed \cite{saha2022tinyml-review} the scope of ML frameworks, tasks, metrics, including a comprehensive review on TinyML stack and deployment pipeline. A number of challenges need to be met in order to fit the tiny resource budgets typical of microcontrollers (kiloBytes of memory, power consumption in mWatt, CPU frequency in MHz) while maintaining performance at an acceptable level and retaining portability to extremely polymorphic hardware in this category.

{\bf Embedded IoT Software Platforms --}
Various open source IoT operating systems are used to provide hardware abstraction, and resource sharing primitives as well as convenient peripheral access (e.g. sensor/actuator, network subsystem) on heterogeneous low-power microcontrollers. Prior work such as \cite{hahm2015operating} surveys such operating systems, among which prominent examples include RIOT\cite{baccelli2018riot} and Zephyr\cite{zephyr}. However, such software platforms so far offer very limited support for ML frameworks -- if at all. Moreover, the most advanced support so far are typically hardware- or vendor-specific e.g. with libraries provided by STM32CubeMx or ARM CMSIS-NN.

{\bf Low-power IoT Testbeds --}
Various testbeds offer remote access to fleets of reprogrammable microcontroller-based devices. Prior work such as \cite{lima2019testbeds,mitton2011testbeds} survey such testbeds, among which prominent examples include the open access facility IoT-lab \cite{adjih2015iot-lab}, which offers remote bare-metal access (serial over TCP) to a fleet composed of hundreds of popular low-power boards of various kind.
 
{\bf Benchmarking Suites for TinyML --}
Benchmarking ML on low-power hardware entails a number of challenges \cite{banbury2020benchmarking}. Prior work such as MLPerf Tiny \cite{banbury2021mlperf} provides a standard benchmark suite (a fixed set of representative ML tasks) for evaluating the performance of given hardware, and an online platform for manufacturers to publish their comparative benchmark results. 
In contrast, U-TOE offers a more powerful and more customizable toolkit for performing feasibility's check of user-defined machines learning models on low-power devices, with a greater degree of flexibility and customization.
 
{\bf TinyML Benchmarks --} 
Prior work such as \cite{osman2022tinyml-bench} focuses on performance comparison of different machine learning frameworks on two COTS low-power boards (Arduino Nano BLE 33 and STM32 NUCLEO-F401RE). In particular, it benchmarked two TinyML frameworks, TFLM and X-CUBE-AI over gesture recognition and wake word spotting. Other work such as \cite{sudharsan2021tinyml-bench} tested TFLM models on several microcontroller-based boards. While such papers provide a performance comparison of specific frameworks on specific boards for specific tasks, U-TOE offers greater flexibility and generality, allowing developers to evaluate a wider range of (user-specified) models on a larger variety of low-power devices, and to dive into the execution details of ML models.
 
{\bf TinyML Model Transpiler \& Compilers --}
Compilers such as TVM \cite{chen2018tvm} can be used to automate the transpilation and compilation of models provided by major ML frameworks (TFML, Pytorch etc.) so as to expose low-level routines and optimize them for execution on specific processing unit characteristics (CPU, GPU etc.). An extension of TVM called uTVM was recently introduced, adding smaller hardware targets including a variety of MCUs (microcontroller units).

{\bf TinyML Model Profilers --}
ML-EXray\cite{qiu2022ml-xray} enables TinyML developers to gain visibility into the layer-level details of ML execution and diagnose cloud-to-edge deployment issues. Developers can analyze and debug edge deployment pipelines with high usability, using less than 15 lines of code for fully examination. However, the reliance on Tensorflow Lite restricts the capability to accommodate models from further ML frameworks and deploy on low-power devices.
Major ML frameworks (TFML, Pytorch and MXNet etc.) provide internal profiler \cite{2023_yousefzadeh-asl-miandoab_ProfilingMonitoringDeep}. Such tools allows developers to measure the performance of their models. They can be used to collect metrics such as inference time and memory usage, which can then be analyzed to optimize the model's performance. Though it can provide us execution details in layer level, it still lacks the supports for on-device deployment and evaluation on various IoT devices, while U-TOE is a more general-purpose toolkit that provides a comprehensive solution on a wide range of low-power devices. The above is summarized in Table \ref{tab:related-work-summary}.

\begin{table}[!ht]
\begin{threeparttable}

    \centering
    \caption{Comparison of AIoT frameworks and toolkits.}
    \label{tab:related-work-summary}
    \begin{tabular}{llllll}
    \hline
        Framework & MCU & \makecell[l]{Model\\Type} & \makecell[l]{Remote\\Eval.} & Granularity & \makecell{Model to\\Board Sol.} \\ \hline
        MLPerf & Yes & \makecell[l]{Specified} & No & Model & No \\ 
        ML-EXray & No & TFLite & No & Layer & No \\ 
        TFLite & Yes & TFLite & No & Layer & No \\ 
        Pytorch & No & Torch & No & Layer & No \\ 
        uTVM  & Yes & Universal & No & Operator & No \\
        \bf U-TOE & Yes & Universal & Yes & Operator & Yes\\ \hline
    \end{tabular}
\begin{tablenotes}
    \item Eval. -- Evaluation, Sol. -- Solution
\end{tablenotes}
\end{threeparttable}
\end{table}